% Separate, double-anonymous supplementary appendix for SIGMOD review.
\documentclass[sigconf,review,anonymous]{acmart}

\setcopyright{none}
\copyrightyear{2027}
\acmYear{2027}
\acmDOI{}
\acmISBN{}
\acmConference[SIGMOD '27]{ACM SIGMOD International Conference on Management of Data}{June 13--19, 2027}{Huntington Beach, CA, USA}
\settopmatter{printacmref=false}

\usepackage{amsmath}
\usepackage{amsthm}
\usepackage{algorithm}
\usepackage{algpseudocode}
\usepackage{booktabs}

\newtheorem{theorem}{Theorem}

\begin{document}

\title{Supplementary Appendix: The Cost of Adaptivity: Online Fanout Selection in Cache-Conscious B-Trees}

\author{Anonymous Author(s)}
\affiliation{%
  \institution{Anonymous Institution}
  \country{}
}

\maketitle

\section{Scope}

This appendix supplies proof detail and experimental parameters for the main
paper.  It does not introduce new claims or experiments, and the main paper is
self-contained.  Reviewers need not read this appendix to evaluate the paper's
contributions.

\section{Break-Even Derivation}

Consider a stable phase containing $T$ operations.  Let $c_0$ be the
average operation cost under the incumbent fanout and $c_1$ the cost under a
candidate fanout, excluding online adaptation overhead.  Define the oracle
benefit per operation as

\[
  \Delta = c_0-c_1.
\]

Let $m$, $\rho$, and $d$ denote the per-operation costs of monitoring,
records metadata, and policy evaluation.  Let $R$ be the one-time cost of
rebuilding and publishing the segment.  Retaining the incumbent costs
$Tc_0$; adapting costs

\[
  Tc_1 + T(m+\rho+d) + R.
\]

Therefore adaptation is profitable exactly when

\[
  T(\Delta-m-\rho-d)>R.
\]

If $\Delta>m+\rho+d$, the minimum profitable phase length is

\[
  T_{\min}=\frac{R}{\Delta-m-\rho-d}.
\]

If $\Delta\le m+\rho+d$, no finite phase length amortizes the recurring
online cost.  This latter case explains the measured in-memory workloads in
the main paper: increasing phase length can amortize rebuilding, but cannot
remove recurring monitoring, records, and policy work.

\section{Potential-Function Detail}

For each segment, let $\Phi_t\ge0$ be adaptation credit immediately before
operation $t$.  Each operation contributes at most a constant $\gamma$.
A rebuild of actual cost $R_t$ is permitted only when $\Phi_t\ge R_t$, and
spends $R_t$ credit.

\begin{theorem}
The rebuild rule gives $O(1)$ amortized restructuring cost per operation.
\end{theorem}

\begin{proof}
Let $a_t\in\{0,R_t\}$ be actual restructuring cost and define

\[
  \widehat a_t=a_t+\Phi_{t+1}-\Phi_t.
\]

Without a rebuild, $a_t=0$ and
$\Phi_{t+1}-\Phi_t\le\gamma$, so
$\widehat a_t\le\gamma$.  With a rebuild,
$a_t=R_t$ and
$\Phi_{t+1}-\Phi_t\le\gamma-R_t$, so again
$\widehat a_t\le\gamma$.  Summing over $T$ operations telescopes the
potential terms:

\[
  \sum_{t=1}^{T}a_t
  \le T\gamma+\Phi_1-\Phi_{T+1}
  \le T\gamma+\Phi_1.
\]

Thus total restructuring cost is $O(T)$, or $O(1)$ amortized per
operation.
\end{proof}

This result bounds restructuring only.  It does not eliminate the recurring
online term $T(m+\rho+d)$, which is why amortized rebuilding alone does not
guarantee an end-to-end speedup.

\section{Oracle Construction}

The phase oracle is intentionally stronger than an implementable online
policy: it knows the phase boundary and hot segment.  For each candidate hot
fanout, it rebuilds only the designated segment and reports both operation-only
throughput and throughput including the measured reconfiguration time.

\begin{algorithm}[t]
\caption{Phase-aware fanout oracle}
\begin{algorithmic}[1]
\Require Phases $P$, candidate fanouts $\mathcal F$, default fanout $f_0$
\For{each candidate hot fanout $f\in\mathcal F$}
  \State initialize all segments with fanout $f_0$
  \For{each phase $p\in P$}
    \State obtain the known hot segment $s_p$
    \State time rebuilding $s_p$ with fanout $f$
    \State replay the fixed operation sequence for $p$
    \State record latency, throughput, and rebuild time
  \EndFor
\EndFor
\State report the best measured candidate per phase and in aggregate
\end{algorithmic}
\end{algorithm}

The oracle is an upper-bound diagnostic rather than a deployable competitor.
If it cannot outperform a static baseline, online adaptation has no fanout
headroom to capture.  If it wins only slightly, the oracle gap bounds the total
monitoring, decision, and rebuilding budget available to an online policy.

\section{Experimental Parameters}

Table~\ref{tab:parameters} records the principal settings used by the paper's
reported experiments.  Workload operation sequences are generated once per
configuration and replayed across index variants so that comparisons use the
same keys and operation types.

\begin{table}[t]
\centering
\caption{Principal benchmark parameters.}
\label{tab:parameters}
\begin{tabular}{ll}
\toprule
Parameter & Setting \\
\midrule
Language standard & C++11 \\
Compiler & GCC 6.3.0 / MinGW \\
Default fanout & 64 \\
Candidate fanouts & 8, 16, 32, 64, 128, 256 \\
Segments & 8 equal key ranges \\
YCSB records & 100,000 \\
YCSB operations & 100,000 per workload \\
Zipf parameter & $0.99$ \\
Adaptation interval & 5,000 operations \\
Repeated trials & 5 \\
Shifting phase length & 30,000 operations \\
Thread benchmark keys & 1,000,000 \\
Thread benchmark operations & 10,000,000 \\
Thread counts & 1, 4, 8, 16, 32 \\
Thread operation mix & 95\% reads, 5\% updates \\
Sliding-hotspot windows & $5\times1{,}000{,}000$ operations \\
Hot range per window & 20\% of key space \\
\bottomrule
\end{tabular}
\end{table}

Repeated-trial summaries report the arithmetic mean, sample standard
deviation, coefficient of variation, and a two-sided 95\% confidence interval
using Student's $t$ distribution with four degrees of freedom.

\section{Reader Safety Boundary}

The concurrent scalability experiment preloads the tree and then executes
lookups and in-place updates.  Existing-value updates do not delete nodes.
Structural splits publish new nodes while retaining old nodes until all worker
threads join, so an optimistic reader can safely validate and retry without
accessing reclaimed memory.  Concurrent segment replacement is outside that
experiment: the adaptive rebuild measurements are single-threaded.  A fully
concurrent adaptive implementation would require epoch-based reclamation (or
an equivalent safe-memory-reclamation mechanism), adding another online cost
to the break-even model.

\section*{Generative AI Disclosure}
Large language models were used to assist with code scaffolding and
boilerplate, and to improve the clarity, organization, and flow of the
manuscript and appendix.  The authors determined the research questions and
system design, conducted and validated the experiments, verified the
mathematical analysis and citations, interpreted the results, and take
responsibility for the work.

\end{document}
